\uuid{C9qo}
\titre{Points stationnaires}
\auteur{hanani ; mzari}
\organisation{exo7}
\licence{CC BY-NC-SA 4.0}
\datecreate{2020-01-20}
\niveau{L1}
\module{Géométrie}
\chapitre{Courbes planes}
\sousChapitre{Courbes paramétrées}
\difficulte{2}
\theme{Courbes paramétrées, points stationnaires}
\exo7qcmid{lille-2.courbes.005}

\contenu{
\begin{question}
\qclassification{244.01}
\qtags{motcle=points stationnaires, branches infinies}
La trajectoire $\Gamma$ d'une particule en mouvement est donnée par les équations
$$x=2t^2-t^4\quad \mbox{et} \quad  y=t^2+t^4.$$
\begin{answers}
\good{Le point de paramètre $0$ est un point de rebroussement de seconde espèce.}
\bad{La tangente à $\Gamma$ au point $(x(0),y(0))$ est la droite d'équation $y=2x$.}
\bad{$\Gamma$ admet la droite d'équation $y=-x$ comme asymptote quand $t$ tend vers l'infini.}
\good{$\Gamma$ admet une branche parabolique de direction asymptotique la droite d'équation $y=-x$.}
\end{answers}
\begin{explanations}
Le $DL_4(0)$ de $f(t)=(x(t),y(t))$ est
$$f(t)=t^2(2,1)+t^4(-1,1)+t^4\left(\varepsilon_1(t),\varepsilon_2(t)\right)\mbox{ avec }\lim _{t\to 0}\varepsilon_i(t)=0.$$
Donc le point de paramètre $0$ est un point de rebroussement de seconde espèce et la tangente en ce point est la droite d'équation $2y-x=0$. Enfin $\displaystyle \lim _{t\to \pm \infty}[y(t)+x(t)]=+\infty$, donc $\Gamma$ admet une branche parabolique de direction asymptotique la droite d'équation $y=-x$.
\end{explanations}
\end{question}
}
